% Part: first-order-logic
% Chapter: syntax-and-semantics
% Section: Semantic Notions

\documentclass[../../../include/open-logic-section]{subfiles}

\begin{document}

\olfileid[pt]{fol}{syn}{sem}
\olsection{Noções Semânticas}

\begin{explain}
Dada a definição de !!{structure}s para linguagens de primeira ordem, podemos
definir algumas propriedades semânticas básicas de sentenças e relações entre
elas. A mais simples destas é a noção de \emph{validade} de
uma sentença. Uma sentença é válida se é satisfeita em toda
!!{structure}. Sentenças válidas são aquelas que são satisfeitas independentemente de
como os símbolos não lógicos nela são interpretados. Sentenças válidas
são, portanto, também chamadas de \emph{verdades lógicas}---elas são verdadeiras, isto é,
satisfeitas, em qualquer !!{structure} e, portanto, sua verdade depende apenas dos
símbolos lógicos que nelas ocorrem e de sua !!{structure} sintática, mas não
dos símbolos não lógicos ou de sua interpretação.
\end{explain}

\begin{defn}[Validade]
Uma sentença $!A$ é \emph{válida}, $\Entails !A$, sse $\Sat{M}{!A}$ para toda
!!{structure}~$\Struct M$.
\end{defn}

\begin{defn}[Consequência]
Uma sentença~$!A$ é \emph{consequência} de um conjunto de sentenças~$\Gamma$, $\Gamma
\Entails !A$, sse para toda !!{structure}~$\Struct M$ com
$\Sat{M}{\Gamma}$, $\Sat{M}{!A}$.
\end{defn}


\begin{defn}[Satisfatibilidade]
Um conjunto de sentenças~$\Gamma$ é \emph{satisfatível} se $\Sat{M}{\Gamma}$
para alguma !!{structure}~$\Struct M$. Se $\Gamma$ não é satisfatível ele é
chamado de \emph{insatisfatível}.
\end{defn}

\begin{prop}
Uma sentença $!A$ é válida sse $\Gamma \Entails !A$ para todo conjunto de
sentenças~$\Gamma$.
\end{prop}

\begin{proof}
Para a direção direta, seja $!A$ válida, e seja $\Gamma$ um
conjunto de sentenças. Seja $\Struct M$ !!a{structure} tal que
$\Sat{M}{\Gamma}$. Como $!A$ é válida, $\Sat{M}{!A}$, portanto $\Gamma
\Entails !A$.

Para a contrapositiva da direção reversa, seja $!A$ inválida,
então há !!a{structure}~$\Struct M$ com $\Sat/{M}{!A}$. Quando $\Gamma
= \{ \ltrue \}$, como $\ltrue$ é válida, $\Sat{M}{\Gamma}$. Portanto,
há !!a{structure}~$\Struct M$ tal que $\Sat{M}{\Gamma}$ mas
$\Sat/{M}{!A}$, portanto $!A$ não é consequência de $\Gamma$.
\end{proof}

\begin{prop}
\ollabel{prop:entails-unsat}
$\Gamma \Entails !A$ sse $\Gamma \cup \{\lnot !A\}$ é insatisfatível.
\end{prop}

\begin{proof}
Para a direção direta, suponha $\Gamma \Entails !A$ e suponha, ao
contrário, que há !!a{structure}~$\Struct M$ tal que $\Sat{M}{\Gamma
  \cup \{ \lnot !A \}}$. Como $\Sat{M}{\Gamma}$ e $\Gamma \Entails
!A$, $\Sat{M}{!A}$. Além disso, como $\Sat{M}{\Gamma\cup \{ \lnot !A \}}$,
$\Sat{M}{\lnot !A}$, então temos tanto $\Sat{M}{!A}$ quanto $\Sat/{M}{!A}$,
uma contradição. Portanto, não pode haver tal !!{structure}~$\Struct M$, então
$\Gamma \cup \{ \lnot !A \}$ é insatisfatível.

Para a direção reversa, suponha que $\Gamma \cup \{ \lnot !A \}$ é
insatisfatível. Então, para toda !!{structure}~$\Struct M$, ou
$\Sat/{M}{\Gamma}$ ou $\Sat{M}{!A}$. Portanto, para toda !!{structure}
$\Struct M$ com $\Sat{M}{\Gamma}$, $\Sat{M}{!A}$, então $\Gamma \Entails
!A$.
\end{proof}

\begin{prob}
\begin{enumerate}
\item Mostre que $\Gamma \Entails \bot$ sse $\Gamma$ é insatisfatível.
\item Mostre que $\Gamma \cup \{!A\} \Entails \bot$ sse $\Gamma \Entails \lnot !A$.
\item Suponha que $c$ não ocorre em $!A$ nem em $\Gamma$. Mostre que
  $\Gamma \Entails \lforall[x][!A]$ sse $\Gamma \Entails
  \Subst{!A}{c}{x}$.
\end{enumerate}
\end{prob}

\begin{prop}
Se $\Gamma \subseteq \Gamma'$ e $\Gamma \Entails !A$, então $\Gamma'
\Entails !A$.
\end{prop}

\begin{proof}
Suponha que $\Gamma \subseteq \Gamma'$ e $\Gamma \Entails !A$. Seja
$\Struct M$ uma estrutura tal que $\Sat{M}{\Gamma'}$; então $\Sat{M}{\Gamma}$,
e como $\Gamma \Entails !A$, obtemos que $\Sat{M}{!A}$. Portanto,
sempre que $\Sat{M}{\Gamma'}$, $\Sat{M}{!A}$, então $\Gamma' \Entails !A$.
\end{proof}


\begin{thm}[Teorema da Dedução Semântica]
\ollabel{thm:sem-deduction}
$\Gamma \cup \{!A\} \Entails !B$ sse $\Gamma \Entails !A \lif !B$.
\end{thm}

\begin{proof}
Para a direção direta, seja $\Gamma \cup \{ !A \} \Entails !B$ e
seja $\Struct M$ !!a{structure} tal que $\Sat{M}{\Gamma}$. Se
$\Sat{M}{!A}$, então $\Sat{M}{\Gamma \cup \{ !A \} }$, então, como $!B$ é consequência de $\Gamma
\cup \{ !A \}$, obtemos $\Sat{M}{!B}$. Portanto,
$\Sat{M}{!A \lif !B}$, então $\Gamma \Entails !A \lif !B$.

Para a direção reversa, seja $\Gamma \Entails !A \lif !B$ e
$\Struct M$ !!a{structure} tal que $\Sat{M}{\Gamma \cup \{ !A
  \}}$. Então $\Sat{M}{\Gamma}$, então $\Sat{M}{!A \lif !B}$, e como
$\Sat{M}{!A}$, $\Sat{M}{!B}$. Portanto, sempre que $\Sat{M}{\Gamma \cup \{
  !A \} }$, $\Sat{M}{!B}$, então $\Gamma \cup \{ !A \} \Entails !B$.
\end{proof}

\begin{prop}\ollabel{prop:quant-terms}
  Seja $\Struct{M}$ !!a{structure}, e $!A(x)$ !!a{formula} com
  uma variável livre~$x$, e $t$~um termo fechado. Então:
  \begin{tagenumerate}{prvEx,prvAll}
    \tagitem{prvEx}{$!A(t) \Entails \lexists[x][!A(x)]$}{}
    \tagitem{prvAll}{$\lforall[x][!A(x)] \Entails !A(t)$}{}
  \end{tagenumerate}
\end{prop}

\begin{proof}
  \begin{tagenumerate}{prvEx,prvAll}
    \tagitem{prvEx}{%
      \iftag{probEx}{Exercício.}{Suponha $\Sat{M}{!A(t)}$. Seja $s$ uma
        atribuição de variáveis com $s(x) = \Value{t}{M}$. Então
        $\Sat{M}{!A(t)}[s]$ já que $!A(t)$ é !!a{sentence}. Por
        \olref[ext]{prop:ext-formulas}, $\Sat{M}{!A(x)}[s]$. Por
        \olref[ass]{prop:sat-quant},
        $\Sat{M}{\lexists[x][!A(x)]}$.}}{}
    \tagitem{prvAll}{%
      \iftag{probAll}{Exercício.}{Suponha
        $\Sat{M}{\lforall[x][!A(x)]}$. Seja $s$ uma atribuição de
        variáveis com $s(x) = \Value{t}{M}$. Por
        \olref[ass]{prop:sat-quant}, $\Sat{M}{!A(x)}[s]$.  Por
        \olref[ext]{prop:ext-formulas}, $\Sat{M}{!A(t)}[s]$. Por
        \olref[ass]{prop:sentence-sat-true}, $\Sat{M}{!A(t)}$ já que
        $!A(t)$ é !!a{sentence}.}}{}
  \end{tagenumerate}
\end{proof}

\begin{probtag}{probEx,probAll}
  Complete a prova de \olref[fol][syn][sem]{prop:quant-terms}.
\end{probtag} 

\end{document}
